\section{Conclusion and Future Work}
\label{sec:conclusion}

We have analyzed the network structure of \module{Mathlib} at three levels of granularity plus their cross-level interaction. Table~\ref{tab:summary} collects the key quantitative findings.

\begin{table}[ht]
\centering
\caption{Cross-level summary of structural findings.}
\label{tab:summary}
\small
\begin{tabular}{lrrrr}
\toprule
Metric & \S\ref{sec:module-import} Module & \S\ref{sec:theorem-premise} Declaration & \S\ref{sec:namespace-graph} Namespace & \S\ref{sec:module-declaration} Cross-Level \\
\midrule
Nodes                      & $7{,}563$        & $308{,}129$      & $10{,}097$       & --- \\
Edges                      & $23{,}570$       & $8{,}436{,}366$  & $332{,}081$      & --- \\
Cross-namespace edges       & $37.1\%$         & $50.9\%$         & $85.8\%$         & --- \\
Intra-module edges          & ---              & $9.6\%$          & ---              & --- \\
Centrality separation       & partial          & complete         & complete         & --- \\
Community--namespace NMI   & ---              & $0.34$           & $0.25$           & --- \\
Modularity                 & ---              & $0.48$           & $0.28$           & --- \\
Single-node removal impact & $29$ components  & $\le 6$ nodes    & GCC $20.8\%$ at $20\%$ & --- \\
Containment (depth $1$)    & ---              & ---              & $22.2\%$         & --- \\
Theorem/lemma in-degree ratio & ---            & $1.47\times$     & ---              & --- \\
Module cohesion (mean)     & ---              & ---              & ---              & $0.107$ \\
Zero-cohesion modules      & ---              & ---              & ---              & $8.4\%$ \\
$G_{\mathrm{file}}/G_{\mathrm{module}}$ edge ratio & ---  & ---  & ---              & $9.1\times$ \\
Active imports             & ---              & ---              & ---              & $72.1\%$ \\
Indirect declaration deps  & ---              & ---              & ---              & $92.2\%$ \\
\bottomrule
\end{tabular}
\end{table}

Rather than recapitulating per-level findings (detailed in Appendix~\ref{app:supplementary}), we organise the discussion around the product\,/\,process lens of \S\ref{sec:conceptual-framework}.

\label{sec:summary}
\label{sec:interplay}
The central claim of this paper is that the network structure of a formal mathematical library records not only the inherited mathematical structure, but also the human organizational process. Our data substantiate this claim with three process-traces embedded in the product, corresponding to the findings of \S\ref{sec:contribution}:

\begin{enumerate}
\item \textbf{Organizational divergence and cognitive compression} (Finding~1). At every granularity level, naming conventions and dependency-based community structure diverge---contributors impose a coarser taxonomy than the dependency graph warrants, trading structural precision for navigability. The module graph compresses the declaration-level network by $9.1\times$ in edges (\S\ref{sec:import_vs_declaration}), with $92.2\%$ of cross-file declaration dependencies reached through transitive import chains. A further asymmetry separates product from process: in-degree is heavy-tailed and open-ended (a product property), while out-degree is bounded by cognitive complexity (a process property).

\item \textbf{Structural redundancy as development ergonomics} (Finding~2). The $17.5\%$ post-\texttt{shake} redundancy (\S\ref{sec:transitive-reduction}) records the gap between compiler requirements and developer workflow. The redundancy is directional---it inflates out-degree without affecting in-degree---confirming that developers import for ergonomics, not along minimal dependency paths.

\item \textbf{Topological inversion and semantic flattening} (Finding~3). The depth inversion---$153$/$84$/$8$ layers across module, declaration, and namespace graphs---records that the file hierarchy is shaped by incremental layering, while mathematical logic resists deep stratification. Formalization preserves traditional labels but flattens their meaning: the \texttt{theorem}/\texttt{lemma} distinction is partially validated in aggregate but contradicted in many individual cases (\S\ref{sec:thm-vs-lemma}).
\end{enumerate}

The relationship between product and process is not unidirectional. The $37.1\%$ cross-directory edge rate in $G_{\mathrm{module}}^{-}$ (\S\ref{sec:cross-namespace}) exerts reorganization pressure on the directory hierarchy---as witnessed by \module{Mathlib}'s ongoing absorption of \module{RingTheory} into \module{Algebra}. The deep dependency chain ($153$ layers) constrains the development process directly: modifications to high-PageRank modules trigger recompilation cascades that incentivize stability at the core and experimentation at the periphery. Formalization thus initiates a feedback cycle: the module tree shapes the module graph; the module graph reshapes the module tree.

\label{sec:practice}
These findings are not merely descriptive; several metrics admit direct operational use. Module cohesion (Definition~\ref{def:cohesion}) can flag scope drift when computed incrementally; the zero-citation rate (\S\ref{sec:thm-leaves}) identifies potentially redundant API surface. PageRank and betweenness rankings identify modules whose modification triggers the widest recompilation cascades, and could inform CI prioritization---assigning higher test priority to pull requests that touch high-centrality modules. Huch and Wenzel~\cite{huch2024} demonstrate ${>}100\times$ speedup from DAG-based parallel builds; Baanen et al.~\cite{baanen2025} report $6$--$33\%$ speedups from structural refactors.

For language design, the two-layer hub structure (\S\ref{sec:thm-types}) suggests that an explicit boundary between language infrastructure and mathematical content could enable independent dependency analysis and compilation. The $92.2\%$ indirect dependency rate (\S\ref{sec:import_vs_declaration}) motivates finer-grained import mechanisms; Lean's module system (November 2025), distinguishing \texttt{public import} from \texttt{import}, is already a step in this direction.

For AI systems that learn premise selection or proof generation from \module{Mathlib}, our analysis reveals structural features usable as training metadata: theorem/lemma annotations as importance signals, Louvain community labels as disciplinary tags, and declaration types as role indicators. As AI-generated proofs become a larger fraction of contributions, the metrics established here---redundancy rate, cohesion, degree distribution exponents, cross-namespace density---provide a baseline for detecting structural drift.

\label{sec:limitations}
Several limitations constrain these conclusions. All statistics are computed from a single commit (\texttt{534cf0b}, February 2026); we cannot determine whether the structural properties we observe are stable, growing, or approaching critical thresholds. The evolutionary hypotheses advanced in \S\ref{sec:cross-level-summary} remain untested predictions. Our snapshot further captures the library in a transitional state where most imports are still public; as the community adopts private imports through Lean's module system, redundancy rates and transitive visibility will change. The ``process'' side of our analysis is inferred from structural traces in the product, not observed directly---we do not analyze PR workflows, contributor patterns, or commit history. Finally, our data are consistent with bidirectional influence between the module tree and the dependency graph (\S\ref{sec:summary}), but a causal account would require longitudinal or experimental evidence.

\label{sec:future}
Several graph definitions in \S\ref{sec:graph-definitions} remain empirically unexplored; \S\ref{sec:temporal-graphs} outlines a programme for temporal, visibility, and co-modification graphs. Beyond these, we highlight three open questions. The growing ecosystem of downstream projects (FLT~\cite{flt_lean}, PFR~\cite{pfr_lean}, Sphere Eversion~\cite{sphere_eversion}, etc.) suggests a ``core + satellite'' architecture whose structural properties merit investigation; the federated-tiers hypothesis (\S\ref{app:module-analysis}) can be reframed as Mathlib consolidating as the foundational tier. Applying the same analysis to other formal libraries---Isabelle/AFP~\cite{blanchette2015,klein2012}, Coq/MathComp, Mizar/MML~\cite{bancerek2018}---would distinguish structural features intrinsic to mathematics from those contingent on a particular proof assistant or community. A deeper question underlies all of these: is the dependency topology we observe a property of \emph{mathematics}, or of \emph{how mathematics is currently produced}? The metrics developed in this paper provide the vocabulary in which the answer, when it comes, can be expressed.
